\documentclass[11pt, ukrainian]{article}
\setlength\textwidth{180mm} \setlength\textheight{277mm}
\setlength\voffset{-38mm} \setlength\hoffset{-28mm}
%\setlength\parindent{0mm}
\pagestyle{empty}
\usepackage[T2A]{fontenc}
\usepackage[utf8]{inputenc}
\usepackage[ukrainian]{babel}
\begin{document}
{\Large \center  Варіанти завдань до лабораторної роботи №4\par 2020/21 н.р.\par \bigskip\bigskip}

Порядок полів у записах та інформацію додаткового файлу фіксовано у варіанті.

Для порівняння числових значень використовується традиційний порядок.

Для порівняння рядкових значень використовується лексикографічний порядок. 

У випадку сортування за кількома полями результуюче сортування будується 
за порядками окремих полів за допомогою конструкції лексикографічного порядку. 
Послідовність полів, за якими має відбуватися сортування, задається в умові. 
Якщо не вказано інше, то поле сортується за зростанням.

Якщо умова вимагає знайти $k$ най… і $k$-е місце займають кілька, то слід виводити всіх їх.

\newpage
\par \bigskip

{\center Варіант  \No \textbf { 1 }\par}
\bigskip
%type = meteo
%variant = 101
Обробляються дані спостережень за погодою. 

\underline{Записи основного файлу містять поля}: мінімальна денна температура, максимальна денна температура, день, місяць, сила вітру, код метеостанції, кількість опадів, середня денна температура, рік, відносна вологість.

\underline{У допоміжному файлі наявні ключі}: найменше значення відносної вологості, Failed successfullY.

\underline{Знайти} дати, по яких середня максимальна денна температура була більша ніж в середньому по всіх спостереженнях. Вивести по кожному з них інформацію:\par
{\noindent}--- на першому рядку:\par
рік, місяць, день, кількість опадів за цей рік, середня середньоденна температура (беремо середні по датах і потім їх середнє);

{\noindent}--- на наступних рядках, починаючи з табуляції, вивести для дати спостереження метеостанції, по яких різниця максимальної та мінімальної температур менша 5.13 (по одному на рядок):\par
середня денна температура, максимальна денна температура, мінімальна денна температура, метеостанція

{\noindent}у такому сортуванні: різниця максимальної та мінімальної температур, метеостанція.

%Дати сортувати: рік, місяць, день.

\bigskip\textit{Обмеження}

Код метеостанції є рядком довжини від 4 до 6 включно.
Може містити літери алфавіту, десяткові цифри, підкреслення. Решта полів числові.

Рік є чотиризначним цілим.
Кількість опадів є дійсним значенням, подається у форматі з 2 знаками після крапки.
Сила вітру є дійсним значенням, подається у форматі з 1 знаком після крапки.
Температури є дійсними значеннями, подаються у форматі з 3 знаками після крапки.
Цілі значення подаються без десяткової крапки.
Для дійсних використовується зображення з фіксованою крапкою.

\textit{Вихідне форматування дійсних}

Усі дійсні значення виводяться у форматі з фіксованою крапкою з тією ж кількістю знаків після крапки, що і у вхідному зображенні.

Якщо у варіанті раніше не було зазначено, то \underline{обчислювані програмою} середні значення виводяться у форматі з 2 знаками після крапки.
%\bigskip%\textit{Для деяких варіантів може знадобитись}


\par
\newpage
{\center Варіант  \No \textbf { 2 }\par}
\bigskip
%type = meteo
%variant = 113
Обробляються дані спостережень за погодою. 

\underline{Записи основного файлу містять поля}: рік, середня денна температура, код метеостанції, місяць, день, відносна вологість, кількість опадів, максимальна денна температура, мінімальна денна температура, сила вітру.

\underline{У допоміжному файлі наявні ключі}: останній рік, зазначений у спостереженнях; кількість метеостанцій, по яких велися спостереження в останньому році.

\underline{Знайти} метеостанції, по яких фіксувалася сила вітру у два рази більша за середню по метеостанції. Вивести по кожному з них інформацію:\par
{\noindent}--- на першому рядку:\par
метеостанція, найбільша сила вітру, середня різниця максимальної та мінімальної денних температур;

{\noindent}--- на наступних рядках, починаючи з табуляції, вивести для них дані по датах (по одному на рядок):\par
середня денна температура, сила вітру, вологість, місяць, день, рік

{\noindent}у такому сортуванні: місяць, день, рік.

%Метеостанції сортувати: найбільша сила вітру (за спаданням), метеостанція.

\bigskip\textit{Обмеження}

Код метеостанції є рядком довжини від 3 до 10 включно.
Може містити літери алфавіту, десяткові цифри, підкреслення. Решта полів числові.

Рік є чотиризначним цілим.
Кількість опадів є цілим значенням.
Сила вітру є дійсним значенням, подається у форматі з 2 знаками після крапки.
Температури є дійсними значеннями, подаються у форматі з 2 знаками після крапки.
Цілі значення подаються без десяткової крапки.
Для дійсних використовується зображення з фіксованою крапкою.

\textit{Вихідне форматування дійсних}

Усі дійсні значення виводяться у форматі з фіксованою крапкою з тією ж кількістю знаків після крапки, що і у вхідному зображенні.

Якщо у варіанті раніше не було зазначено, то \underline{обчислювані програмою} середні значення виводяться у форматі з 3 знаками після крапки.
%\bigskip%\textit{Для деяких варіантів може знадобитись}


\par
\newpage
{\center Варіант  \No \textbf { 3 }\par}
\bigskip
%type = meteo
%variant = 96
Обробляються дані спостережень за погодою. 

\underline{Записи основного файлу містять поля}: рік, середня денна температура, мінімальна денна температура, відносна вологість, місяць, день, кількість опадів, код метеостанції, максимальна денна температура, сила вітру.

\underline{У допоміжному файлі наявні ключі}: найбільша зафіксована кількість опадів, загальна кількість дат, коли фіксувалися спостереження.

\underline{Знайти} дати, по яких фіксувалася максимальна денна температура менша ніж в середньому по всіх спостереженнях. Вивести по кожному з них інформацію:\par
{\noindent}--- на першому рядку:\par
день, місяць, рік, найбільша та найменша вологість за дату, середня кількість опадів;

{\noindent}--- на наступних рядках, починаючи з табуляції, вивести для дати спостереження метеостанції, де вологість була більша 60 за відсутності опадів (по одному на рядок):\par
вологість, середня температура, метеостанція

{\noindent}у такому сортуванні: вологість, метеостанція.

%Дати сортувати: місяць, середня кількість опадів, рік, день.

\bigskip\textit{Обмеження}

Код метеостанції є рядком довжини від 5 до 11 включно.
Може містити літери алфавіту, десяткові цифри, підкреслення. Решта полів числові.

Рік є чотиризначним цілим.
Кількість опадів є дійсним значенням, подається у форматі з 1 знаком після крапки.
Сила вітру є цілим значенням.
Температури є дійсними значеннями, подаються у форматі з 1 знаком після крапки.
Цілі значення подаються без десяткової крапки.
Для дійсних використовується зображення з фіксованою крапкою.

\textit{Вихідне форматування дійсних}

Усі дійсні значення виводяться у форматі з фіксованою крапкою з тією ж кількістю знаків після крапки, що і у вхідному зображенні.

Якщо у варіанті раніше не було зазначено, то \underline{обчислювані програмою} середні значення виводяться у форматі з 1 знаком після крапки.
%\bigskip%\textit{Для деяких варіантів може знадобитись}


\par
\newpage
{\center Варіант  \No \textbf { 4 }\par}
\bigskip
%type = meteo
%variant = 108
Обробляються дані спостережень за погодою. 

\underline{Записи основного файлу містять поля}: сила вітру, відносна вологість, середня денна температура, кількість опадів, мінімальна денна температура, максимальна денна температура, день, місяць, рік, код метеостанції.

\underline{У допоміжному файлі наявні ключі}: найбільша середня температура, найменша вологість.

\underline{Знайти} 8 метеостанцій, на яких була зафіксована найменша максимальна денна температура. Вивести по кожному з них інформацію:\par
{\noindent}--- на першому рядку:\par
середні значення показників: середня, максимальна та мінімальна денна  температура, вологість, сила вітру; та код метеостанції;

{\noindent}--- на наступних рядках, починаючи з табуляції, вивести для них дані по датах (по одному на рядок):\par
день, рік, місяць, різниця максимальної та мінімальної температур, вологість

{\noindent}у такому сортуванні: вологість, місяць, день, рік.

%Метеостанція сортувати: середня вологість, середня сила вітру, метеостанція.

\bigskip\textit{Обмеження}

Код метеостанції є рядком довжини від 4 до 8 включно.
Може містити літери алфавіту, десяткові цифри, підкреслення. Решта полів числові.

Рік є чотиризначним цілим.
Кількість опадів є дійсним значенням, подається у форматі з 2 знаками після крапки.
Сила вітру є дійсним значенням, подається у форматі з 1 знаком після крапки.
Температури є дійсними значеннями, подаються у форматі з 1 знаком після крапки.
Цілі значення подаються без десяткової крапки.
Для дійсних використовується зображення з фіксованою крапкою.

\textit{Вихідне форматування дійсних}

Усі дійсні значення виводяться у форматі з фіксованою крапкою з тією ж кількістю знаків після крапки, що і у вхідному зображенні.

Якщо у варіанті раніше не було зазначено, то \underline{обчислювані програмою} середні значення виводяться у форматі з 3 знаками після крапки.
%\bigskip%\textit{Для деяких варіантів може знадобитись}


\par
\newpage
{\center Варіант  \No \textbf { 5 }\par}
\bigskip
%type = meteo
%variant = 94
Обробляються дані спостережень за погодою. 

\underline{Записи основного файлу містять поля}: максимальна денна температура, середня денна температура, мінімальна денна температура, рік, сила вітру, кількість опадів, код метеостанції, відносна вологість, день, місяць.

\underline{У допоміжному файлі наявні ключі}: найбільша максимальна температура, найменша максимальна температура, SMILE.

\underline{Знайти} дати, по яких фіксувалася кількість опадів, більша ніж в середньому по всіх спостереженнях. Вивести по кожному з них інформацію:\par
{\noindent}--- на першому рядку:\par
рік, місяць, день, найбільша кількість опадів за цю дату, середня сила вітру;

{\noindent}--- на наступних рядках, починаючи з табуляції, вивести для дати спостереження метеостанції, коди яких починаються з 13 (по одному на рядок):\par
метеостанція, вологість, середня денна температура

{\noindent}у такому сортуванні: метеостанція.

%Дати сортувати: середня сила вітру, кількість опадів, рік, місяць, день.

\bigskip\textit{Обмеження}

Код метеостанції є рядком довжини від 5 до 7 включно.
Може містити літери алфавіту, десяткові цифри, підкреслення. Решта полів числові.

Рік є чотиризначним цілим.
Кількість опадів є цілим значенням.
Сила вітру є дійсним значенням, подається у форматі з 2 знаками після крапки.
Температури є дійсними значеннями, подаються у форматі з 3 знаками після крапки.
Цілі значення подаються без десяткової крапки.
Для дійсних використовується зображення з фіксованою крапкою.

\textit{Вихідне форматування дійсних}

Усі дійсні значення виводяться у форматі з фіксованою крапкою з тією ж кількістю знаків після крапки, що і у вхідному зображенні.

Якщо у варіанті раніше не було зазначено, то \underline{обчислювані програмою} середні значення виводяться у форматі з 2 знаками після крапки.
%\bigskip%\textit{Для деяких варіантів може знадобитись}


\par
\newpage
{\center Варіант  \No \textbf { 6 }\par}
\bigskip
%type = meteo
%variant = 97
Обробляються дані спостережень за погодою. 

\underline{Записи основного файлу містять поля}: день, місяць, відносна вологість, рік, кількість опадів, максимальна денна температура, код метеостанції, середня денна температура, сила вітру, мінімальна денна температура.

\underline{У допоміжному файлі наявні ключі}: середня сила вітру, сума всіх місяців.

\underline{Знайти} дати, по яких фіксувалася мінімальна денна температура менше ніж в середньому по всіх спостереженнях. Вивести по кожному з них інформацію:\par
{\noindent}--- на першому рядку:\par
день, місяць, рік, кількість опадів у цьому місяці, максимальна сила вітру у цьому місяці;

{\noindent}--- на наступних рядках, починаючи з табуляції, вивести для дати спостереження метеостанції, по яких в цю дату не було опадів (по одному на рядок):\par
середня денна температура, максимальна денна температура, мінімальна денна температура, сила вітру, метеостанція

{\noindent}у такому сортуванні: середня денна температура, метеостанція.

%Дати сортувати: місяць, день, рік.

\bigskip\textit{Обмеження}

Код метеостанції є рядком довжини від 3 до 12 включно.
Може містити літери алфавіту, десяткові цифри, підкреслення. Решта полів числові.

Рік є чотиризначним цілим.
Кількість опадів є дійсним значенням, подається у форматі з 1 знаком після крапки.
Сила вітру є цілим значенням.
Температури є дійсними значеннями, подаються у форматі з 2 знаками після крапки.
Цілі значення подаються без десяткової крапки.
Для дійсних використовується зображення з фіксованою крапкою.

\textit{Вихідне форматування дійсних}

Усі дійсні значення виводяться у форматі з фіксованою крапкою з тією ж кількістю знаків після крапки, що і у вхідному зображенні.

Якщо у варіанті раніше не було зазначено, то \underline{обчислювані програмою} середні значення виводяться у форматі з 1 знаком після крапки.
%\bigskip%\textit{Для деяких варіантів може знадобитись}


\par
\newpage
{\center Варіант  \No \textbf { 7 }\par}
\bigskip
%type = meteo
%variant = 83
Обробляються дані спостережень за погодою. 

\underline{Записи основного файлу містять поля}: місяць, день, кількість опадів, відносна вологість, код метеостанції, середня денна температура, мінімальна денна температура, максимальна денна температура, сила вітру, рік.

\underline{У допоміжному файлі наявні ключі}: кількість різних дат, довжина найдовшого коду метеостанції.

\underline{Знайти} 6 дат, за які випала найменша кількість опадів (сумарно по всіх метеостанціях). Вивести по кожному з них інформацію:\par
{\noindent}--- на першому рядку:\par
кількість опадів, рік, місяць, день, середня вологість;

{\noindent}--- на наступних рядках, починаючи з табуляції, вивести для дати спостереження метеостанцій, де сила вітру була більша за 3.2 (по одному на рядок):\par
середня денна температура, максимальна денна температура, мінімальна денна температура, сила вітру, метеостанція

{\noindent}у такому сортуванні: середня денна температура, метеостанція.

%Дати сортувати: місяць, день, кількість опадів, рік.

\bigskip\textit{Обмеження}

Код метеостанції є рядком довжини від 5 до 9 включно.
Може містити літери алфавіту, десяткові цифри, підкреслення. Решта полів числові.

Рік є чотиризначним цілим.
Кількість опадів є цілим значенням.
Сила вітру є цілим значенням.
Температури є дійсними значеннями, подаються у форматі з 1 знаком після крапки.
Цілі значення подаються без десяткової крапки.
Для дійсних використовується зображення з фіксованою крапкою.

\textit{Вихідне форматування дійсних}

Усі дійсні значення виводяться у форматі з фіксованою крапкою з тією ж кількістю знаків після крапки, що і у вхідному зображенні.

Якщо у варіанті раніше не було зазначено, то \underline{обчислювані програмою} середні значення виводяться у форматі з 2 знаками після крапки.
%\bigskip%\textit{Для деяких варіантів може знадобитись}


\par
\newpage
{\center Варіант  \No \textbf { 8 }\par}
\bigskip
%type = meteo
%variant = 100
Обробляються дані спостережень за погодою. 

\underline{Записи основного файлу містять поля}: місяць, день, сила вітру, середня денна температура, код метеостанції, кількість опадів, максимальна денна температура, мінімальна денна температура, рік, відносна вологість.

\underline{У допоміжному файлі наявні ключі}: кількість днів, в які велися спостереження; загальна кількість записів.

\underline{Знайти} дати, по яких найбільша середньоденна температура була більша ніж в середньому по всіх спостереженнях. Вивести по кожному з них інформацію:\par
{\noindent}--- на першому рядку:\par
день, рік, місяць, кількість опадів, середня середньоденна температура, найбільша середньоденна температура;

{\noindent}--- на наступних рядках, починаючи з табуляції, вивести для дати спостереження метеостанції, по яких денна температура відрізняється від середньої середньоденної принаймні на 10 (по одному на рядок):\par
мінімальна денна температура, максимальна денна температура, середня денна температура, вологість, метеостанція

{\noindent}у такому сортуванні: різниця максимальної та мінімальної температур, метеостанція.

%Дати сортувати: кількість опадів, місяць, рік, день.

\bigskip\textit{Обмеження}

Код метеостанції є рядком довжини від 4 до 8 включно.
Може містити літери алфавіту, десяткові цифри, підкреслення. Решта полів числові.

Рік є чотиризначним цілим.
Кількість опадів є дійсним значенням, подається у форматі з 1 знаком після крапки.
Сила вітру є дійсним значенням, подається у форматі з 1 знаком після крапки.
Температури є дійсними значеннями, подаються у форматі з 2 знаками після крапки.
Цілі значення подаються без десяткової крапки.
Для дійсних використовується зображення з фіксованою крапкою.

\textit{Вихідне форматування дійсних}

Усі дійсні значення виводяться у форматі з фіксованою крапкою з тією ж кількістю знаків після крапки, що і у вхідному зображенні.

Якщо у варіанті раніше не було зазначено, то \underline{обчислювані програмою} середні значення виводяться у форматі з 3 знаками після крапки.
%\bigskip%\textit{Для деяких варіантів може знадобитись}


\par
\newpage
{\center Варіант  \No \textbf { 9 }\par}
\bigskip
%type = meteo
%variant = 87
Обробляються дані спостережень за погодою. 

\underline{Записи основного файлу містять поля}: кількість опадів, рік, відносна вологість, максимальна денна температура, сила вітру, код метеостанції, день, місяць, мінімальна денна температура, середня денна температура.

\underline{У допоміжному файлі наявні ключі}: сумарна кількість опадів, найменша мінімальна температура.

\underline{Знайти} 5 дат, по яких фіксувалася найменша середня денна температура. Вивести по кожному з них інформацію:\par
{\noindent}--- на першому рядку:\par
рік, місяць, день, кількість опадів (сумарна), найменша середня денна температура;

{\noindent}--- на наступних рядках, починаючи з табуляції, вивести для дати спостереження метеостанції, де середня денна температура була не більш ніж 10 більша за найменшу середню денну температуру за відповідну дату (по одному на рядок):\par
метеостанція, сила вітру, вологість

{\noindent}у такому сортуванні: метеостанція (за спаданням).

%Дати сортувати: день, місяць, кількість опадів, день.

\bigskip\textit{Обмеження}

Код метеостанції є рядком довжини від 3 до 10 включно.
Може містити літери алфавіту, десяткові цифри, підкреслення. Решта полів числові.

Рік є чотиризначним цілим.
Кількість опадів є дійсним значенням, подається у форматі з 2 знаками після крапки.
Сила вітру є дійсним значенням, подається у форматі з 2 знаками після крапки.
Температури є дійсними значеннями, подаються у форматі з 3 знаками після крапки.
Цілі значення подаються без десяткової крапки.
Для дійсних використовується зображення з фіксованою крапкою.

\textit{Вихідне форматування дійсних}

Усі дійсні значення виводяться у форматі з фіксованою крапкою з тією ж кількістю знаків після крапки, що і у вхідному зображенні.

Якщо у варіанті раніше не було зазначено, то \underline{обчислювані програмою} середні значення виводяться у форматі з 1 знаком після крапки.
%\bigskip%\textit{Для деяких варіантів може знадобитись}


\par
\newpage
{\center Варіант  \No \textbf { 10 }\par}
\bigskip
%type = meteo
%variant = 111
Обробляються дані спостережень за погодою. 

\underline{Записи основного файлу містять поля}: середня денна температура, кількість опадів, відносна вологість, сила вітру, мінімальна денна температура, максимальна денна температура, код метеостанції, місяць, день, рік.

\underline{У допоміжному файлі наявні ключі}: перший рік, в якому були спостереження; кількість спостережень, в яких середня денна температура була більша 10.

\underline{Знайти} 3 метеостанції, на яких була зафіксована найбільша кількість опадів за день. Вивести по кожному з них інформацію:\par
{\noindent}--- на першому рядку:\par
кількість спостережень, середня кількість опадів, найбільша кількість опадів за день, метеостанція;

{\noindent}--- на наступних рядках, починаючи з табуляції, вивести для них дані по місяцях (по одному на рядок):\par
рік, місяць, середня сила вітру за місяць, кількість опадів за місяць

{\noindent}у такому сортуванні: місяць, рік.

%Метеостанції сортувати: кількість спостережень (за спаданням), метеостанція.

\bigskip\textit{Обмеження}

Код метеостанції є рядком довжини від 3 до 6 включно.
Може містити літери алфавіту, десяткові цифри, підкреслення. Решта полів числові.

Рік є чотиризначним цілим.
Кількість опадів є дійсним значенням, подається у форматі з 1 знаком після крапки.
Сила вітру є дійсним значенням, подається у форматі з 1 знаком після крапки.
Температури є дійсними значеннями, подаються у форматі з 3 знаками після крапки.
Цілі значення подаються без десяткової крапки.
Для дійсних використовується зображення з фіксованою крапкою.

\textit{Вихідне форматування дійсних}

Усі дійсні значення виводяться у форматі з фіксованою крапкою з тією ж кількістю знаків після крапки, що і у вхідному зображенні.

Якщо у варіанті раніше не було зазначено, то \underline{обчислювані програмою} середні значення виводяться у форматі з 1 знаком після крапки.
%\bigskip%\textit{Для деяких варіантів може знадобитись}


\par
\newpage
{\center Варіант  \No \textbf { 11 }\par}
\bigskip
%type = meteo
%variant = 107
Обробляються дані спостережень за погодою. 

\underline{Записи основного файлу містять поля}: кількість опадів, відносна вологість, рік, код метеостанції, місяць, день, середня денна температура, максимальна денна температура, мінімальна денна температура, сила вітру.

\underline{У допоміжному файлі наявні ключі}: найбільша кількість опадів у спостереженнях, найменша кількість опадів у спостереженнях, загальна кількість записів.

\underline{Знайти} 4 метеостанції, на яких була зафіксована найбільша середньоденна температура. Вивести по кожному з них інформацію:\par
{\noindent}--- на першому рядку:\par
середня вологість, середня середньоденна температура, найбільша середньоденна температура, найбільша сила вітру, метеостанція;

{\noindent}--- на наступних рядках, починаючи з табуляції, вивести для них дані по датах (по одному на рядок):\par
день, місяць, день, вологість, середня денна температура, максимальна денна температура, сила вітру, вологість, рік

{\noindent}у такому сортуванні: вологість (за спаданням), місяць, день, рік.

%Метеостанція сортувати: кількість спостережень (за спаданням), найбільша сила вітру, метеостанція.

\bigskip\textit{Обмеження}

Код метеостанції є рядком довжини від 4 до 11 включно.
Може містити літери алфавіту, десяткові цифри, підкреслення. Решта полів числові.

Рік є чотиризначним цілим.
Кількість опадів є цілим значенням.
Сила вітру є дійсним значенням, подається у форматі з 2 знаками після крапки.
Температури є дійсними значеннями, подаються у форматі з 2 знаками після крапки.
Цілі значення подаються без десяткової крапки.
Для дійсних використовується зображення з фіксованою крапкою.

\textit{Вихідне форматування дійсних}

Усі дійсні значення виводяться у форматі з фіксованою крапкою з тією ж кількістю знаків після крапки, що і у вхідному зображенні.

Якщо у варіанті раніше не було зазначено, то \underline{обчислювані програмою} середні значення виводяться у форматі з 3 знаками після крапки.
%\bigskip%\textit{Для деяких варіантів може знадобитись}


\par
\newpage
{\center Варіант  \No \textbf { 12 }\par}
\bigskip
%type = meteo
%variant = 88
Обробляються дані спостережень за погодою. 

\underline{Записи основного файлу містять поля}: рік, сила вітру, мінімальна денна температура, кількість опадів, день, місяць, відносна вологість, середня денна температура, код метеостанції, максимальна денна температура.

\underline{У допоміжному файлі наявні ключі}: середня вологість по всіх спостереженнях, загальна кількість спостережень.

\underline{Знайти} 6 дат, по яких фіксувалася найбільша максимальна денна температура. Вивести по кожному з них інформацію:\par
{\noindent}--- на першому рядку:\par
день, рік, місяць, відношення найбільшої максимальної температури до найменшої мінімальної (по всіх метеостанціях, округлене до 2 знаків);

{\noindent}--- на наступних рядках, починаючи з табуляції, вивести для дати спостереження метеостанції, де вологість не перевищувала 45 (по одному на рядок):\par
вологість, метеостанція, сила вітру, середня денна температура

{\noindent}у такому сортуванні: середня денна температура, метеостанція.

%Дати сортувати: відношення (округлене), рік, місяць, день.

\bigskip\textit{Обмеження}

Код метеостанції є рядком довжини від 5 до 12 включно.
Може містити літери алфавіту, десяткові цифри, підкреслення. Решта полів числові.

Рік є чотиризначним цілим.
Кількість опадів є дійсним значенням, подається у форматі з 2 знаками після крапки.
Сила вітру є цілим значенням.
Температури є дійсними значеннями, подаються у форматі з 1 знаком після крапки.
Цілі значення подаються без десяткової крапки.
Для дійсних використовується зображення з фіксованою крапкою.

\textit{Вихідне форматування дійсних}

Усі дійсні значення виводяться у форматі з фіксованою крапкою з тією ж кількістю знаків після крапки, що і у вхідному зображенні.

Якщо у варіанті раніше не було зазначено, то \underline{обчислювані програмою} середні значення виводяться у форматі з 2 знаками після крапки.
%\bigskip%\textit{Для деяких варіантів може знадобитись}


\par
\newpage
{\center Варіант  \No \textbf { 13 }\par}
\bigskip
%type = meteo
%variant = 104
Обробляються дані спостережень за погодою. 

\underline{Записи основного файлу містять поля}: кількість опадів, рік, мінімальна денна температура, код метеостанції, максимальна денна температура, місяць, день, сила вітру, відносна вологість, середня денна температура.

\underline{У допоміжному файлі наявні ключі}: найбільша різниця денних температур у спостереженнях, кількість записів.

\underline{Знайти} 4 метеостанції, по яких у середньому була найвища середньоденна температура. Вивести по кожному з них інформацію:\par
{\noindent}--- на першому рядку:\par
кількість спостережень, середня середньоденна температура найменша вологість, метеостанція;

{\noindent}--- на наступних рядках, починаючи з табуляції, вивести для них дані по датах, за які не було опадів (по одному на рядок):\par
середня денна температура, максимальна денна температура, мінімальна денна температура, сила вітру, вологість, місяць, день, рік

{\noindent}у такому сортуванні: вологість, рік, місяць, день.

%Метеостанції сортувати: найменша вологість, кількість спостережень (за спаданням), метеостанція.

\bigskip\textit{Обмеження}

Код метеостанції є рядком довжини від 5 до 9 включно.
Може містити літери алфавіту, десяткові цифри, підкреслення. Решта полів числові.

Рік є чотиризначним цілим.
Кількість опадів є цілим значенням.
Сила вітру є дійсним значенням, подається у форматі з 1 знаком після крапки.
Температури є дійсними значеннями, подаються у форматі з 2 знаками після крапки.
Цілі значення подаються без десяткової крапки.
Для дійсних використовується зображення з фіксованою крапкою.

\textit{Вихідне форматування дійсних}

Усі дійсні значення виводяться у форматі з фіксованою крапкою з тією ж кількістю знаків після крапки, що і у вхідному зображенні.

Якщо у варіанті раніше не було зазначено, то \underline{обчислювані програмою} середні значення виводяться у форматі з 2 знаками після крапки.
%\bigskip%\textit{Для деяких варіантів може знадобитись}


\par
\newpage
{\center Варіант  \No \textbf { 14 }\par}
\bigskip
%type = meteo
%variant = 93
Обробляються дані спостережень за погодою. 

\underline{Записи основного файлу містять поля}: рік, сила вітру, середня денна температура, код метеостанції, місяць, день, кількість опадів, відносна вологість, максимальна денна температура, мінімальна денна температура.

\underline{У допоміжному файлі наявні ключі}: кількість спостережень за вересень, загальна кількість записів.

\underline{Знайти} дати, по яких фіксувалася відносна вологість більша ніж в середньому по всіх спостереженнях. Вивести по кожному з них інформацію:\par
{\noindent}--- на першому рядку:\par
найбільше та найменше значення вологості за дату, середня середньоденна температура, рік, місяць, день, кількість опадів;

{\noindent}--- на наступних рядках, починаючи з табуляції, вивести для дати спостереження метеостанції, де фіксувалися опади (по одному на рядок):\par
метеостанція, кількість опадів, вологість

{\noindent}у такому сортуванні: кількість опадів (за спаданням), метеостанція.

%Дати сортувати: найбільша вологість (за спаданням), рік, місяць, день.

\bigskip\textit{Обмеження}

Код метеостанції є рядком довжини від 3 до 7 включно.
Може містити літери алфавіту, десяткові цифри, підкреслення. Решта полів числові.

Рік є чотиризначним цілим.
Кількість опадів є дійсним значенням, подається у форматі з 2 знаками після крапки.
Сила вітру є дійсним значенням, подається у форматі з 2 знаками після крапки.
Температури є дійсними значеннями, подаються у форматі з 3 знаками після крапки.
Цілі значення подаються без десяткової крапки.
Для дійсних використовується зображення з фіксованою крапкою.

\textit{Вихідне форматування дійсних}

Усі дійсні значення виводяться у форматі з фіксованою крапкою з тією ж кількістю знаків після крапки, що і у вхідному зображенні.

Якщо у варіанті раніше не було зазначено, то \underline{обчислювані програмою} середні значення виводяться у форматі з 1 знаком після крапки.
%\bigskip%\textit{Для деяких варіантів може знадобитись}


\par
\newpage
{\center Варіант  \No \textbf { 15 }\par}
\bigskip
%type = meteo
%variant = 103
Обробляються дані спостережень за погодою. 

\underline{Записи основного файлу містять поля}: максимальна денна температура, код метеостанції, кількість опадів, мінімальна денна температура, середня денна температура, відносна вологість, сила вітру, рік, день, місяць.

\underline{У допоміжному файлі наявні ключі}: середня сила вітру, середня відносна вологість.

\underline{Знайти} 5 метеостанцій, по яких у середньому випала найбільша кількість опадів. Вивести по кожному з них інформацію:\par
{\noindent}--- на першому рядку:\par
метеостанція, кількість спостережень, середня кількість опадів, найбільша сила вітру;

{\noindent}--- на наступних рядках, починаючи з табуляції, вивести для них дані по датах (по одному на рядок):\par
рік, місяць, день, середня денна температура, максимальна денна температура, мінімальна денна температура, сила вітру, вологість

{\noindent}у такому сортуванні: місяць, день, рік.

%Метеостанції сортувати: кількість спостережень (за спаданням), найбільша сила вітру, метеостанція.

\bigskip\textit{Обмеження}

Код метеостанції є рядком довжини від 4 до 11 включно.
Може містити літери алфавіту, десяткові цифри, підкреслення. Решта полів числові.

Рік є чотиризначним цілим.
Кількість опадів є дійсним значенням, подається у форматі з 1 знаком після крапки.
Сила вітру є цілим значенням.
Температури є дійсними значеннями, подаються у форматі з 1 знаком після крапки.
Цілі значення подаються без десяткової крапки.
Для дійсних використовується зображення з фіксованою крапкою.

\textit{Вихідне форматування дійсних}

Усі дійсні значення виводяться у форматі з фіксованою крапкою з тією ж кількістю знаків після крапки, що і у вхідному зображенні.

Якщо у варіанті раніше не було зазначено, то \underline{обчислювані програмою} середні значення виводяться у форматі з 3 знаками після крапки.
%\bigskip%\textit{Для деяких варіантів може знадобитись}


\par
\newpage
{\center Варіант  \No \textbf { 16 }\par}
\bigskip
%type = meteo
%variant = 89
Обробляються дані спостережень за погодою. 

\underline{Записи основного файлу містять поля}: рік, кількість опадів, середня денна температура, місяць, день, відносна вологість, максимальна денна температура, сила вітру, мінімальна денна температура, код метеостанції.

\underline{У допоміжному файлі наявні ключі}: рік завершення спостережень, місяць завершення спостережень, середня сила вітру.

\underline{Знайти} 8 дат, по яких фіксувалася найменша максимальна  денна температура. Вивести по кожному з них інформацію:\par
{\noindent}--- на першому рядку:\par
відношення кількості опадів (середньої за дату) до вологості (середньої за дату) (округлене до 3 знаків), місяць, день, рік, середня максимальна температура;

{\noindent}--- на наступних рядках, починаючи з табуляції, вивести для дати спостереження метеостанції, де сила вітру була принаймні тричі більша за вологість (по одному на рядок):\par
середня денна температура, максимальна денна температура, мінімальна денна температура, сила вітру, вологість, метеостанція

{\noindent}у такому сортуванні: сила вітру, метеостанція.

%Дати сортувати: відношення (округлене, за спаданням), місяць, день, рік.

\bigskip\textit{Обмеження}

Код метеостанції є рядком довжини від 3 до 8 включно.
Може містити літери алфавіту, десяткові цифри, підкреслення. Решта полів числові.

Рік є чотиризначним цілим.
Кількість опадів є дійсним значенням, подається у форматі з 1 знаком після крапки.
Сила вітру є цілим значенням.
Температури є дійсними значеннями, подаються у форматі з 3 знаками після крапки.
Цілі значення подаються без десяткової крапки.
Для дійсних використовується зображення з фіксованою крапкою.

\textit{Вихідне форматування дійсних}

Усі дійсні значення виводяться у форматі з фіксованою крапкою з тією ж кількістю знаків після крапки, що і у вхідному зображенні.

Якщо у варіанті раніше не було зазначено, то \underline{обчислювані програмою} середні значення виводяться у форматі з 1 знаком після крапки.
%\bigskip%\textit{Для деяких варіантів може знадобитись}


\par
\newpage
{\center Варіант  \No \textbf { 17 }\par}
\bigskip
%type = meteo
%variant = 109
Обробляються дані спостережень за погодою. 

\underline{Записи основного файлу містять поля}: код метеостанції, максимальна денна температура, мініма\-ль\-на денна температура, середня денна температура, кількість опадів, відносна вологість, сила вітру, рік, місяць, день.

\underline{У допоміжному файлі наявні ключі}: найбільша сила вітру, найменша сила вітру, загальна кількість спостережень.

\underline{Знайти} 5 метеостанцій, по яких у середньому випала найбільша кількість опадів. Вивести по кожному з них інформацію:\par
{\noindent}--- на першому рядку:\par
метеостанція, кількість спостережень, середня кількість опадів, сумарна кількість опадів;

{\noindent}--- на наступних рядках, починаючи з табуляції, вивести для них дані по датах, коли не було опадів (по одному на рядок):\par
середня денна температура, сила вітру, вологість, день, місяць, рік

{\noindent}у такому сортуванні: вологість (за спаданням), рік, місяць, день.

%Метеостанції сортувати: кількість спостережень (за спаданням), сумарна кількість опадів,  метеостанція.

\bigskip\textit{Обмеження}

Код метеостанції є рядком довжини від 4 до 7 включно.
Може містити літери алфавіту, десяткові цифри, підкреслення. Решта полів числові.

Рік є чотиризначним цілим.
Кількість опадів є дійсним значенням, подається у форматі з 2 знаками після крапки.
Сила вітру є дійсним значенням, подається у форматі з 1 знаком після крапки.
Температури є дійсними значеннями, подаються у форматі з 2 знаками після крапки.
Цілі значення подаються без десяткової крапки.
Для дійсних використовується зображення з фіксованою крапкою.

\textit{Вихідне форматування дійсних}

Усі дійсні значення виводяться у форматі з фіксованою крапкою з тією ж кількістю знаків після крапки, що і у вхідному зображенні.

Якщо у варіанті раніше не було зазначено, то \underline{обчислювані програмою} середні значення виводяться у форматі з 2 знаками після крапки.
%\bigskip%\textit{Для деяких варіантів може знадобитись}


\par
\newpage
{\center Варіант  \No \textbf { 18 }\par}
\bigskip
%type = meteo
%variant = 86
Обробляються дані спостережень за погодою. 

\underline{Записи основного файлу містять поля}: код метеостанції, середня денна температура, кількість опадів, рік, відносна вологість, сила вітру, місяць, день, максимальна денна температура, мінімальна денна температура.

\underline{У допоміжному файлі наявні ключі}: загальна кількість записів, сумарна сила вітру.

\underline{Знайти} 8 дат, по яких фіксувалася найбільша середня денна температура. Вивести по кожному з них інформацію:\par
{\noindent}--- на першому рядку:\par
середня середньоденна температура, найбільша середня денна температура, місяць, рік, день, кількість опадів;

{\noindent}--- на наступних рядках, починаючи з табуляції, вивести для дати спостереження метеостанції, де сила вітру була меншою за середню по даті (по одному на рядок):\par
вологість, сила вітру, середня температура, метеостанція

{\noindent}у такому сортуванні: вологість, метеостанція.

%Дати сортувати: середня середньоденна температура (за спаданням), рік, місяць, день.

\bigskip\textit{Обмеження}

Код метеостанції є рядком довжини від 5 до 6 включно.
Може містити літери алфавіту, десяткові цифри, підкреслення. Решта полів числові.

Рік є чотиризначним цілим.
Кількість опадів є цілим значенням.
Сила вітру є дійсним значенням, подається у форматі з 2 знаками після крапки.
Температури є дійсними значеннями, подаються у форматі з 1 знаком після крапки.
Цілі значення подаються без десяткової крапки.
Для дійсних використовується зображення з фіксованою крапкою.

\textit{Вихідне форматування дійсних}

Усі дійсні значення виводяться у форматі з фіксованою крапкою з тією ж кількістю знаків після крапки, що і у вхідному зображенні.

Якщо у варіанті раніше не було зазначено, то \underline{обчислювані програмою} середні значення виводяться у форматі з 3 знаками після крапки.
%\bigskip%\textit{Для деяких варіантів може знадобитись}


\par
\newpage
{\center Варіант  \No \textbf { 19 }\par}
\bigskip
%type = meteo
%variant = 85
Обробляються дані спостережень за погодою. 

\underline{Записи основного файлу містять поля}: код метеостанції, середня денна температура, рік, відносна вологість, мінімальна денна температура, сила вітру, місяць, день, кількість опадів, максимальна денна температура.

\underline{У допоміжному файлі наявні ключі}: рік початку спостережень, кількість записів по 13-их числах.

\underline{Знайти} 7 дат, по яких фіксувалася найбільша сила вітру. Вивести по кожному з них інформацію:\par
{\noindent}--- на першому рядку:\par
найбільша сила вітру, найменша сила вітру, сумарна кількість опадів, місяць, рік, день, середня вологість;

{\noindent}--- на наступних рядках, починаючи з табуляції, вивести для дати спостереження метеостанції, де вологість була менша за середню вологість за дату (по одному на рядок):\par
вологість, метеостанція, сила вітру

{\noindent}у такому сортуванні: вологість, сила вітру (за спаданням), метеостанція.

%Дати сортувати: рік, місяць, день, сумарна кількість опадів.

\bigskip\textit{Обмеження}

Код метеостанції є рядком довжини від 5 до 12 включно.
Може містити літери алфавіту, десяткові цифри, підкреслення. Решта полів числові.

Рік є чотиризначним цілим.
Кількість опадів є дійсним значенням, подається у форматі з 2 знаками після крапки.
Сила вітру є цілим значенням.
Температури є дійсними значеннями, подаються у форматі з 2 знаками після крапки.
Цілі значення подаються без десяткової крапки.
Для дійсних використовується зображення з фіксованою крапкою.

\textit{Вихідне форматування дійсних}

Усі дійсні значення виводяться у форматі з фіксованою крапкою з тією ж кількістю знаків після крапки, що і у вхідному зображенні.

Якщо у варіанті раніше не було зазначено, то \underline{обчислювані програмою} середні значення виводяться у форматі з 2 знаками після крапки.
%\bigskip%\textit{Для деяких варіантів може знадобитись}


\par
\newpage
{\center Варіант  \No \textbf { 20 }\par}
\bigskip
%type = meteo
%variant = 99
Обробляються дані спостережень за погодою. 

\underline{Записи основного файлу містять поля}: мінімальна денна температура, відносна вологість, сила вітру, середня денна температура, максимальна денна температура, код метеостанції, місяць, день, кількість опадів, рік.

\underline{У допоміжному файлі наявні ключі}: найменша мінімальна температура, найбільша сила вітру.

\underline{Знайти} дати, по яких середня кількість опадів була більша ніж в середньому по всіх спостереженнях. Вивести по кожному з них інформацію:\par
{\noindent}--- на першому рядку:\par
рік, місяць, день, кількість опадів, середня вологість;

{\noindent}--- на наступних рядках, починаючи з табуляції, вивести для дати спостереження метеостанції, де вологість була більша 50 (по одному на рядок):\par
сила вітру, вологість, метеостанція

{\noindent}у такому сортуванні: вологість, метеостанція.

%Дати сортувати: середня вологість, день, місяць, рік.

\bigskip\textit{Обмеження}

Код метеостанції є рядком довжини від 4 до 9 включно.
Може містити літери алфавіту, десяткові цифри, підкреслення. Решта полів числові.

Рік є чотиризначним цілим.
Кількість опадів є цілим значенням.
Сила вітру є дійсним значенням, подається у форматі з 1 знаком після крапки.
Температури є дійсними значеннями, подаються у форматі з 3 знаками після крапки.
Цілі значення подаються без десяткової крапки.
Для дійсних використовується зображення з фіксованою крапкою.

\textit{Вихідне форматування дійсних}

Усі дійсні значення виводяться у форматі з фіксованою крапкою з тією ж кількістю знаків після крапки, що і у вхідному зображенні.

Якщо у варіанті раніше не було зазначено, то \underline{обчислювані програмою} середні значення виводяться у форматі з 1 знаком після крапки.
%\bigskip%\textit{Для деяких варіантів може знадобитись}


\par
\newpage
{\center Варіант  \No \textbf { 21 }\par}
\bigskip
%type =     meteo
%variant = 82
Обробляються дані спостережень за погодою. 

\underline{Записи основного файлу містять поля}: максимальна денна температура, середня денна температура, день, місяць, мінімальна денна температура, кількість опадів, сила вітру, відносна вологість, рік, код метеостанції.

\underline{У допоміжному файлі наявні ключі}: кількість записів, найбільша сила вітру.

\underline{Знайти} 5 дат, за які випала найбільша кількість опадів (сумарно по всіх метеостанціях). Вивести по кожному з них інформацію:\par
{\noindent}--- на першому рядку:\par
рік, місяць, день, кількість опадів, середня середньоденна температура;

{\noindent}--- на наступних рядках, починаючи з табуляції, вивести для дати спостереження метеостанцій, де максимальна денна температура була більша за середню (по всіх метеостанціях) середньоденну за цю дату (по одному на рядок):\par
метеостанція, середня денна температура, максимальна денна температура, мінімальна денна температура, сила вітру

{\noindent}у такому сортуванні: сила вітру, метеостанція.

%Дати сортувати: кількість опадів, середня середньоденна температура, рік, місяць, день.

\bigskip\textit{Обмеження}

Код метеостанції є рядком довжини від 3 до 10 включно.
Може містити літери алфавіту, десяткові цифри, підкреслення. Решта полів числові.

Рік є чотиризначним цілим.
Кількість опадів є дійсним значенням, подається у форматі з 1 знаком після крапки.
Сила вітру є дійсним значенням, подається у форматі з 2 знаками після крапки.
Температури є дійсними значеннями, подаються у форматі з 1 знаком після крапки.
Цілі значення подаються без десяткової крапки.
Для дійсних використовується зображення з фіксованою крапкою.

\textit{Вихідне форматування дійсних}

Усі дійсні значення виводяться у форматі з фіксованою крапкою з тією ж кількістю знаків після крапки, що і у вхідному зображенні.

Якщо у варіанті раніше не було зазначено, то \underline{обчислювані програмою} середні значення виводяться у форматі з 3 знаками після крапки.
%\bigskip%\textit{Для деяких варіантів може знадобитись}


\par
\newpage
{\center Варіант  \No \textbf { 22 }\par}
\bigskip
%type = meteo
%variant = 90
Обробляються дані спостережень за погодою. 

\underline{Записи основного файлу містять поля}: рік, максимальна денна температура, день, місяць, кількість опадів, код метеостанції, мінімальна денна температура, відносна вологість, сила вітру, середня денна температура.

\underline{У допоміжному файлі наявні ключі}: кількість спостережень, коли сила вітру перевищувала 15; загальна кількість записів.

\underline{Знайти} 7 дат, по яких фіксувалася найбільша мінімальна денна температура. Вивести по кожному з них інформацію:\par
{\noindent}--- на першому рядку:\par
рік, місяць, день, мінімальна температура, середня температура, максимальна температура, вологість;

{\noindent}--- на наступних рядках, починаючи з табуляції, вивести для дати спостереження метеостанції, де сила вітру була меша середньої за дату (по одному на рядок):\par
сила вітру, метеостанція, вологість, різниця максимальної та мінімальної денних температур

{\noindent}у такому сортуванні: різниця максимальної та мінімальної температур, метеостанція.

%Дати сортувати: вологість, місяць, день, рік.

\bigskip\textit{Обмеження}

Код метеостанції є рядком довжини від 4 до 7 включно.
Може містити літери алфавіту, десяткові цифри, підкреслення. Решта полів числові.

Рік є чотиризначним цілим.
Кількість опадів є дійсним значенням, подається у форматі з 2 знаками після крапки.
Сила вітру є цілим значенням.
Температури є дійсними значеннями, подаються у форматі з 2 знаками після крапки.
Цілі значення подаються без десяткової крапки.
Для дійсних використовується зображення з фіксованою крапкою.

\textit{Вихідне форматування дійсних}

Усі дійсні значення виводяться у форматі з фіксованою крапкою з тією ж кількістю знаків після крапки, що і у вхідному зображенні.

Якщо у варіанті раніше не було зазначено, то \underline{обчислювані програмою} середні значення виводяться у форматі з 2 знаками після крапки.
%\bigskip%\textit{Для деяких варіантів може знадобитись}


\par
\newpage
{\center Варіант  \No \textbf { 23 }\par}
\bigskip
%type = meteo
%variant = 95
Обробляються дані спостережень за погодою. 

\underline{Записи основного файлу містять поля}: відносна вологість, код метеостанції, рік, кількість опадів, середня денна температура, максимальна денна температура, місяць, день, мінімальна денна температура, сила вітру.

\underline{У допоміжному файлі наявні ключі}: максимальна сила вітру, що спостерігалася; загальна кількість записів.

\underline{Знайти} дати, по яких фіксувалася середньоденна температура більша ніж в середньому по всіх спостереженнях. Вивести по кожному з них інформацію:\par
{\noindent}--- на першому рядку:\par
рік, місяць, день, кількість опадів за місяць дати;

{\noindent}--- на наступних рядках, починаючи з табуляції, вивести для дати спостереження метеостанції, по яких спостерігалися опади більше 5 (по одному на рядок):\par
метеостанція, кількість опадів, сила вітру

{\noindent}у такому сортуванні: кількість опадів, метеостанція.

%Дати сортувати: рік, місяць, день.

\bigskip\textit{Обмеження}

Код метеостанції є рядком довжини від 3 до 10 включно.
Може містити літери алфавіту, десяткові цифри, підкреслення. Решта полів числові.

Рік є чотиризначним цілим.
Кількість опадів є дійсним значенням, подається у форматі з 1 знаком після крапки.
Сила вітру є дійсним значенням, подається у форматі з 2 знаками після крапки.
Температури є дійсними значеннями, подаються у форматі з 1 знаком після крапки.
Цілі значення подаються без десяткової крапки.
Для дійсних використовується зображення з фіксованою крапкою.

\textit{Вихідне форматування дійсних}

Усі дійсні значення виводяться у форматі з фіксованою крапкою з тією ж кількістю знаків після крапки, що і у вхідному зображенні.

Якщо у варіанті раніше не було зазначено, то \underline{обчислювані програмою} середні значення виводяться у форматі з 3 знаками після крапки.
%\bigskip%\textit{Для деяких варіантів може знадобитись}


\par
\newpage
{\center Варіант  \No \textbf { 24 }\par}
\bigskip
%type = meteo
%variant = 102
Обробляються дані спостережень за погодою. 

\underline{Записи основного файлу містять поля}: сила вітру, місяць, день, кількість опадів, максимальна денна температура, код метеостанції, середня денна температура, рік, мінімальна денна температура, відносна вологість.

\underline{У допоміжному файлі наявні ключі}: кількість спостережень, останній рік спостережень.

\underline{Знайти} дати, по яких середня мінімальна денна температура була менше ніж в середньому по всіх спостереженнях. Вивести по кожному з них інформацію:\par
{\noindent}--- на першому рядку:\par
місяць, день, середня кількість опадів, середня мінімальна температура, рік;

{\noindent}--- на наступних рядках, починаючи з табуляції, вивести для дати спостереження метеостанції, по яких в цей день опадів було менше 3 (по одному на рядок):\par
середня денна температура, метеостанція, сила вітру

{\noindent}у такому сортуванні: середня денна температура, метеостанція.

%Дати сортувати: рік, середня кількість опадів, місяць, день.

\bigskip\textit{Обмеження}

Код метеостанції є рядком довжини від 5 до 9 включно.
Може містити літери алфавіту, десяткові цифри, підкреслення. Решта полів числові.

Рік є чотиризначним цілим.
Кількість опадів є цілим значенням.
Сила вітру є дійсним значенням, подається у форматі з 1 знаком після крапки.
Температури є дійсними значеннями, подаються у форматі з 3 знаками після крапки.
Цілі значення подаються без десяткової крапки.
Для дійсних використовується зображення з фіксованою крапкою.

\textit{Вихідне форматування дійсних}

Усі дійсні значення виводяться у форматі з фіксованою крапкою з тією ж кількістю знаків після крапки, що і у вхідному зображенні.

Якщо у варіанті раніше не було зазначено, то \underline{обчислювані програмою} середні значення виводяться у форматі з 1 знаком після крапки.
%\bigskip%\textit{Для деяких варіантів може знадобитись}


\par
\newpage
{\center Варіант  \No \textbf { 25 }\par}
\bigskip
%type = meteo
%variant = 98
Обробляються дані спостережень за погодою. 

\underline{Записи основного файлу містять поля}: код метеостанції, сила вітру, рік, максимальна денна температура, відносна вологість, кількість опадів, середня денна температура, мінімальна денна температура, день, місяць.

\underline{У допоміжному файлі наявні ключі}: загальна кількість записів, перший рік спостережень.

\underline{Знайти} дати, по яких фіксувалась відносна вологість більша ніж в середньому по всіх спостереженнях. Вивести по кожному з них інформацію:\par
{\noindent}--- на першому рядку:\par
кількість опадів, середня середньоденна температура, місяць, день, рік;

{\noindent}--- на наступних рядках, починаючи з табуляції, вивести для дати спостереження метеостанції, коди яких закінчуються на 17 (по одному на рядок):\par
метеостанція, вологість, середня денна температура, максимальна денна температура, мінімальна денна температура, сила вітру

{\noindent}у такому сортуванні: метеостанція.

%Дати сортувати: середня середньоденна температура, рік (за спаданням), місяць, день.

\bigskip\textit{Обмеження}

Код метеостанції є рядком довжини від 5 до 12 включно.
Може містити літери алфавіту, десяткові цифри, підкреслення. Решта полів числові.

Рік є чотиризначним цілим.
Кількість опадів є дійсним значенням, подається у форматі з 1 знаком після крапки.
Сила вітру є дійсним значенням, подається у форматі з 1 знаком після крапки.
Температури є дійсними значеннями, подаються у форматі з 2 знаками після крапки.
Цілі значення подаються без десяткової крапки.
Для дійсних використовується зображення з фіксованою крапкою.

\textit{Вихідне форматування дійсних}

Усі дійсні значення виводяться у форматі з фіксованою крапкою з тією ж кількістю знаків після крапки, що і у вхідному зображенні.

Якщо у варіанті раніше не було зазначено, то \underline{обчислювані програмою} середні значення виводяться у форматі з 2 знаками після крапки.
%\bigskip%\textit{Для деяких варіантів може знадобитись}


\par
\newpage
{\center Варіант  \No \textbf { 26 }\par}
\bigskip
%type = meteo
%variant = 84
Обробляються дані спостережень за погодою. 

\underline{Записи основного файлу містять поля}: відносна вологість, день, місяць, код метеостанції, сила вітру, кількість опадів, рік, максимальна денна температура, середня денна температура, мінімальна денна температура.

\underline{У допоміжному файлі наявні ключі}: кількість різних дат, кількість записів.

\underline{Знайти} 4 дати, по яких фіксувалася найменша сила вітру. Вивести по кожному з них інформацію:\par
{\noindent}--- на першому рядку:\par
місяць, рік, день, сила вітру, сумарна кількість опадів;

{\noindent}--- на наступних рядках, починаючи з табуляції, вивести для дати спостереження метеостанції (по одному на рядок):\par
сила вітру, кількість опадів, середня температура, код метеостанції

{\noindent}у такому сортуванні: середня температура, метеостанція.

%Дати сортувати: день, кількість опадів, рік, місяць.

\bigskip\textit{Обмеження}

Код метеостанції є рядком довжини від 4 до 6 включно.
Може містити літери алфавіту, десяткові цифри, підкреслення. Решта полів числові.

Рік є чотиризначним цілим.
Кількість опадів є дійсним значенням, подається у форматі з 2 знаками після крапки.
Сила вітру є цілим значенням.
Температури є дійсними значеннями, подаються у форматі з 3 знаками після крапки.
Цілі значення подаються без десяткової крапки.
Для дійсних використовується зображення з фіксованою крапкою.

\textit{Вихідне форматування дійсних}

Усі дійсні значення виводяться у форматі з фіксованою крапкою з тією ж кількістю знаків після крапки, що і у вхідному зображенні.

Якщо у варіанті раніше не було зазначено, то \underline{обчислювані програмою} середні значення виводяться у форматі з 1 знаком після крапки.
%\bigskip%\textit{Для деяких варіантів може знадобитись}


\par
\newpage
{\center Варіант  \No \textbf { 27 }\par}
\bigskip
%type = meteo
%variant = 92
Обробляються дані спостережень за погодою. 

\underline{Записи основного файлу містять поля}: відносна вологість, сила вітру, рік, кількість опадів, мінімальна денна температура, день, місяць, код метеостанції, максимальна денна температура, середня денна температура.

\underline{У допоміжному файлі наявні ключі}: кількість спостережень, в яких не спостерігалися опади; кількість спостережень за січень.

\underline{Знайти} дати, по яких фіксувалася сила вітру більша ніж в середньому по всіх спостереженнях. Вивести по кожному з них інформацію:\par
{\noindent}--- на першому рядку:\par
місяць, рік, день, середня сила вітру, найбільша сила вітру, сумарна кількість опадів за дату;

{\noindent}--- на наступних рядках, починаючи з табуляції, вивести для дати спостереження метеостанції, де сила вітру була менша ніж середня по всіх спостереженнях (всіх станцій) (по одному на рядок):\par
сила вітру, вологість, метеостанція

{\noindent}у такому сортуванні: метеостанція.

%Дати сортувати: день, рік, місяць, найбільша сила вітру.

\bigskip\textit{Обмеження}

Код метеостанції є рядком довжини від 3 до 11 включно.
Може містити літери алфавіту, десяткові цифри, підкреслення. Решта полів числові.

Рік є чотиризначним цілим.
Кількість опадів є цілим значенням.
Сила вітру є дійсним значенням, подається у форматі з 2 знаками після крапки.
Температури є дійсними значеннями, подаються у форматі з 1 знаком після крапки.
Цілі значення подаються без десяткової крапки.
Для дійсних використовується зображення з фіксованою крапкою.

\textit{Вихідне форматування дійсних}

Усі дійсні значення виводяться у форматі з фіксованою крапкою з тією ж кількістю знаків після крапки, що і у вхідному зображенні.

Якщо у варіанті раніше не було зазначено, то \underline{обчислювані програмою} середні значення виводяться у форматі з 3 знаками після крапки.
%\bigskip%\textit{Для деяких варіантів може знадобитись}


\par
\newpage
{\center Варіант  \No \textbf { 28 }\par}
\bigskip
%type = meteo
%variant = 115
Обробляються дані спостережень за погодою. 

\underline{Записи основного файлу містять поля}: середня денна температура, місяць, день, кількість опадів, мінімальна денна температура, рік, код метеостанції, відносна вологість, максимальна денна температура, сила вітру.

\underline{У допоміжному файлі наявні ключі}: кількість записів, максимальна зафіксована вологість, ALL IS DONE.

\underline{Знайти} метеостанції, по яких жодний з показників не відхилявся від середніх по метеостанції більше ніж на 40%. Вивести по кожному з них інформацію:\par
{\noindent}--- на першому рядку:\par
метеостанція, кількість спостережень, найбільше, найменше та середнє значення вологості;

{\noindent}--- на наступних рядках, починаючи з табуляції, вивести для них дані по датах (по одному на рядок):\par
рік, місяць, день, середня денна температура, максимальна денна температура, мінімальна денна температура, сила вітру, вологість

{\noindent}у такому сортуванні: середня денна температура, місяць, день, рік.

%Метеостанції сортувати: найбільше значення вологості (за спаданням), кількість спостережень, метеостанція.

\bigskip\textit{Обмеження}

Код метеостанції є рядком довжини від 5 до 8 включно.
Може містити літери алфавіту, десяткові цифри, підкреслення. Решта полів числові.

Рік є чотиризначним цілим.
Кількість опадів є цілим значенням.
Сила вітру є дійсним значенням, подається у форматі з 1 знаком після крапки.
Температури є дійсними значеннями, подаються у форматі з 3 знаками після крапки.
Цілі значення подаються без десяткової крапки.
Для дійсних використовується зображення з фіксованою крапкою.

\textit{Вихідне форматування дійсних}

Усі дійсні значення виводяться у форматі з фіксованою крапкою з тією ж кількістю знаків після крапки, що і у вхідному зображенні.

Якщо у варіанті раніше не було зазначено, то \underline{обчислювані програмою} середні значення виводяться у форматі з 3 знаками після крапки.
%\bigskip%\textit{Для деяких варіантів може знадобитись}


\par
\newpage
{\center Варіант  \No \textbf { 29 }\par}
\bigskip
%type = meteo
%variant = 110
Обробляються дані спостережень за погодою. 

\underline{Записи основного файлу містять поля}: мінімальна денна температура, день, місяць, максимальна денна температура, відносна вологість, середня денна температура, код метеостанції, кількість опадів, сила вітру, рік.

\underline{У допоміжному файлі наявні ключі}: середня сила вітру, загальна кількість записів.

\underline{Знайти} 6 метеостанцій, на яких була зафіксована найбільша сила вітру. Вивести по кожному з них інформацію:\par
{\noindent}--- на першому рядку:\par
найбільша сила вітру по метеостанції, найменша вологість, середня кількість опадів, метеостанція;

{\noindent}--- на наступних рядках, починаючи з табуляції, вивести для них дані по датах, коли сила вітру була хоча б 50% від максимальної зафіксованої (по одному на рядок):\par
місяць, день, вологість, середня денна температура, максимальна денна температура, мінімальна денна температура, сила вітру, рік

{\noindent}у такому сортуванні: рік, місяць, день.

%Метеостанції сортувати: сила вітру (за спаданням), метеостанція.

\bigskip\textit{Обмеження}

Код метеостанції є рядком довжини від 3 до 8 включно.
Може містити літери алфавіту, десяткові цифри, підкреслення. Решта полів числові.

Рік є чотиризначним цілим.
Кількість опадів є дійсним значенням, подається у форматі з 2 знаками після крапки.
Сила вітру є цілим значенням.
Температури є дійсними значеннями, подаються у форматі з 1 знаком після крапки.
Цілі значення подаються без десяткової крапки.
Для дійсних використовується зображення з фіксованою крапкою.

\textit{Вихідне форматування дійсних}

Усі дійсні значення виводяться у форматі з фіксованою крапкою з тією ж кількістю знаків після крапки, що і у вхідному зображенні.

Якщо у варіанті раніше не було зазначено, то \underline{обчислювані програмою} середні значення виводяться у форматі з 1 знаком після крапки.
%\bigskip%\textit{Для деяких варіантів може знадобитись}


\par
\newpage
{\center Варіант  \No \textbf { 30 }\par}
\bigskip
%type = meteo
%variant = 91
Обробляються дані спостережень за погодою. 

\underline{Записи основного файлу містять поля}: середня денна температура, рік, сила вітру, код метеостанції, відносна вологість, максимальна денна температура, мінімальна денна температура, кількість опадів, місяць, день.

\underline{У допоміжному файлі наявні ключі}: середня силу вітру, загальна кількість спостережень.

\underline{Знайти} 4 дати, по яких фіксувалася найменша мінімальна  денна температура. Вивести по кожному з них інформацію:\par
{\noindent}--- на першому рядку:\par
найменша мінімальна температура, найбільша мінімальна температура, середня вологість, середня сила вітру, середня кількість опадів, місяць, день, рік;

{\noindent}--- на наступних рядках, починаючи з табуляції, вивести для дати спостереження метеостанції, де сила вітру не перевищувала 12 (по одному на рядок):\par
метеостанція, максимальна денна температура, мінімальна денна температура, сила вітру, вологість

{\noindent}у такому сортуванні: сила вітру (за спаданням), метеостанція.

%Дати сортувати: середня вологість, середня кількість опадів, рік, місяць, день.

\bigskip\textit{Обмеження}

Код метеостанції є рядком довжини від 4 до 12 включно.
Може містити літери алфавіту, десяткові цифри, підкреслення. Решта полів числові.

Рік є чотиризначним цілим.
Кількість опадів є дійсним значенням, подається у форматі з 1 знаком після крапки.
Сила вітру є дійсним значенням, подається у форматі з 2 знаками після крапки.
Температури є дійсними значеннями, подаються у форматі з 2 знаками після крапки.
Цілі значення подаються без десяткової крапки.
Для дійсних використовується зображення з фіксованою крапкою.

\textit{Вихідне форматування дійсних}

Усі дійсні значення виводяться у форматі з фіксованою крапкою з тією ж кількістю знаків після крапки, що і у вхідному зображенні.

Якщо у варіанті раніше не було зазначено, то \underline{обчислювані програмою} середні значення виводяться у форматі з 2 знаками після крапки.
%\bigskip%\textit{Для деяких варіантів може знадобитись}


\par
\newpage
{\center Варіант  \No \textbf { 31 }\par}
\bigskip
%type = meteo
%variant = 105
Обробляються дані спостережень за погодою. 

\underline{Записи основного файлу містять поля}: рік, максимальна денна температура, сила вітру, середня денна температура, мінімальна денна температура, відносна вологість, код метеостанції, день, місяць, кількість опадів.

\underline{У допоміжному файлі наявні ключі}: кількість спостережень, найбільше значення кількості опадів, найменше значення кількості опадів.

\underline{Знайти} 6 метеостанцій, по яких у середньому була найвища максимальна денна температура. Вивести по кожному з них інформацію:\par
{\noindent}--- на першому рядку:\par
метеостанція, найбільша та найменша максимальні температури;

{\noindent}--- на наступних рядках, починаючи з табуляції, вивести для них дані по місяцях (по одному на рядок):\par
рік, місяць, кількість опадів, середня середньоденна температура

{\noindent}у такому сортуванні: місяць, рік.

%Метеостанції сортувати: кількість спостережень, найбільша максимальна температура, метеостанція.

\bigskip\textit{Обмеження}

Код метеостанції є рядком довжини від 3 до 9 включно.
Може містити літери алфавіту, десяткові цифри, підкреслення. Решта полів числові.

Рік є чотиризначним цілим.
Кількість опадів є дійсним значенням, подається у форматі з 1 знаком після крапки.
Сила вітру є цілим значенням.
Температури є дійсними значеннями, подаються у форматі з 1 знаком після крапки.
Цілі значення подаються без десяткової крапки.
Для дійсних використовується зображення з фіксованою крапкою.

\textit{Вихідне форматування дійсних}

Усі дійсні значення виводяться у форматі з фіксованою крапкою з тією ж кількістю знаків після крапки, що і у вхідному зображенні.

Якщо у варіанті раніше не було зазначено, то \underline{обчислювані програмою} середні значення виводяться у форматі з 3 знаками після крапки.
%\bigskip%\textit{Для деяких варіантів може знадобитись}


\par
\newpage
{\center Варіант  \No \textbf { 32 }\par}
\bigskip
%type = meteo
%variant = 112
Обробляються дані спостережень за погодою. 

\underline{Записи основного файлу містять поля}: середня денна температура, кількість опадів, код метеостанції, сила вітру, рік, день, місяць, відносна вологість, максимальна денна температура, мінімальна денна температура.

\underline{У допоміжному файлі наявні ключі}: кількість метеостанцій, середнє значення середньоденної температури.

\underline{Знайти} метеостанції, по яких фіксувалася кількість опадів у два рази менша за середню по всіх даних. Вивести по кожному з них інформацію:\par
{\noindent}--- на першому рядку:\par
найбільша кількість опадів, кількість спостережень, середня мінімальна температура, метеостанція;

{\noindent}--- на наступних рядках, починаючи з табуляції, вивести для них агреговані дані (по одному на рядок):\par
день, середня вологість (за цей день року по всіх місяцях по цій метеостанції), рік

{\noindent}у такому сортуванні: день, рік.

%Метеостанції сортувати: середня мінімальна температура, кількість спостережень (за спаданням), метеостанція.

\bigskip\textit{Обмеження}

Код метеостанції є рядком довжини від 4 до 6 включно.
Може містити літери алфавіту, десяткові цифри, підкреслення. Решта полів числові.

Рік є чотиризначним цілим.
Кількість опадів є дійсним значенням, подається у форматі з 2 знаками після крапки.
Сила вітру є дійсним значенням, подається у форматі з 2 знаками після крапки.
Температури є дійсними значеннями, подаються у форматі з 3 знаками після крапки.
Цілі значення подаються без десяткової крапки.
Для дійсних використовується зображення з фіксованою крапкою.

\textit{Вихідне форматування дійсних}

Усі дійсні значення виводяться у форматі з фіксованою крапкою з тією ж кількістю знаків після крапки, що і у вхідному зображенні.

Якщо у варіанті раніше не було зазначено, то \underline{обчислювані програмою} середні значення виводяться у форматі з 2 знаками після крапки.
%\bigskip%\textit{Для деяких варіантів може знадобитись}


\par
\newpage
{\center Варіант  \No \textbf { 33 }\par}
\bigskip
%type = meteo
%variant = 106
Обробляються дані спостережень за погодою. 

\underline{Записи основного файлу містять поля}: середня денна температура, максимальна денна температура, відносна вологість, рік, мінімальна денна температура, кількість опадів, день, місяць, код метеостанції, сила вітру.

\underline{У допоміжному файлі наявні ключі}: найбільша вологість, найменша вологість, кількість різних метеостанцій.

\underline{Знайти} 7 метеостанцій, по яких у середньому була найвища мінімальна денна температура. Вивести по кожному з них інформацію:\par
{\noindent}--- на першому рядку:\par
метеостанція, середня мінімальна температура, найбільша вологість;

{\noindent}--- на наступних рядках, починаючи з табуляції, вивести для них дані по датах (по одному на рядок):\par
місяць, день, середня денна температура, максимальна денна температура, мінімальна денна температура, сила вітру, вологість, рік

{\noindent}у такому сортуванні: сила вітру, місяць, день, рік.

%Метеостанції сортувати: кількість спостережень (за спаданням), найбільша вологість (за спаданням), метеостанція.

\bigskip\textit{Обмеження}

Код метеостанції є рядком довжини від 5 до 7 включно.
Може містити літери алфавіту, десяткові цифри, підкреслення. Решта полів числові.

Рік є чотиризначним цілим.
Кількість опадів є цілим значенням.
Сила вітру є дійсним значенням, подається у форматі з 1 знаком після крапки.
Температури є дійсними значеннями, подаються у форматі з 2 знаками після крапки.
Цілі значення подаються без десяткової крапки.
Для дійсних використовується зображення з фіксованою крапкою.

\textit{Вихідне форматування дійсних}

Усі дійсні значення виводяться у форматі з фіксованою крапкою з тією ж кількістю знаків після крапки, що і у вхідному зображенні.

Якщо у варіанті раніше не було зазначено, то \underline{обчислювані програмою} середні значення виводяться у форматі з 1 знаком після крапки.
%\bigskip%\textit{Для деяких варіантів може знадобитись}


\par
\newpage
{\center Варіант  \No \textbf { 34 }\par}
\bigskip
%type = meteo
%variant = 114
Обробляються дані спостережень за погодою. 

\underline{Записи основного файлу містять поля}: день, місяць, код метеостанції, максимальна денна температура, рік, мінімальна денна температура, середня денна температура, кількість опадів, сила вітру, відносна вологість.

\underline{У допоміжному файлі наявні ключі}: середнє значення дня, округлене до 1 знаку; загальна кількість записів.

\underline{Знайти} метеостанції, по яких жодний з показників не відхилявся від середніх по всіх спостереженнях більше ніж на 50%. Вивести по кожному з них інформацію:\par
{\noindent}--- на першому рядку:\par
метеостанція, середні показники по метеостанції: кількість опадів, сила вітру, середня денна температура;

{\noindent}--- на наступних рядках, починаючи з табуляції, вивести для них дані по місяцях (по одному на рядок):\par
рік, місяць, середня середньоденна температура, середня вологість

{\noindent}у такому сортуванні: рік, місяць.

%Метеостанції сортувати: кількість спостережень (за спаданням), середня сила вітру, метеостанція.

\bigskip\textit{Обмеження}

Код метеостанції є рядком довжини від 5 до 11 включно.
Може містити літери алфавіту, десяткові цифри, підкреслення. Решта полів числові.

Рік є чотиризначним цілим.
Кількість опадів є дійсним значенням, подається у форматі з 1 знаком після крапки.
Сила вітру є цілим значенням.
Температури є дійсними значеннями, подаються у форматі з 1 знаком після крапки.
Цілі значення подаються без десяткової крапки.
Для дійсних використовується зображення з фіксованою крапкою.

\textit{Вихідне форматування дійсних}

Усі дійсні значення виводяться у форматі з фіксованою крапкою з тією ж кількістю знаків після крапки, що і у вхідному зображенні.

Якщо у варіанті раніше не було зазначено, то \underline{обчислювані програмою} середні значення виводяться у форматі з 2 знаками після крапки.
%\bigskip%\textit{Для деяких варіантів може знадобитись}


\par
\newpage
\end{document}
